\documentclass[12pt]{article}
\usepackage{latexblog}

% ============================================================
% Blog Metadata
% ============================================================
\blogtitle{Compile live555 for Android}
\blogdate{2017-11-20}
\blogtags{Android, 移动端, 闲话编程}
\bloglang{zh}
\blogsummary{}

\begin{document}
\makeblogtitle

编译 live555的库在 android 上使用。

首先下载liveMedia库。解压完成可以先在linux环境下编译一遍试试，例如：

\begin{Shaded}
\begin{Highlighting}[]
\ExtensionTok{./genMakefiles}\NormalTok{ macosx}
\FunctionTok{make}
\end{Highlighting}
\end{Shaded}

然后，可以利用ndk-build将它交叉编译成动态库。这时候，需要新建一个Android.mk和Application.mk文件： * Application.mk

\begin{Shaded}
\begin{Highlighting}[]
\DataTypeTok{APP\_BUILD\_SCRIPT} \CharTok{:=}\StringTok{ Android.mk}
\DataTypeTok{APP\_STL} \CharTok{:=}\StringTok{ gnustl\_shared}
\DataTypeTok{APP\_ABI} \CharTok{:=}\StringTok{ armeabi{-}v7a}
\end{Highlighting}
\end{Shaded}

\begin{itemize}
\tightlist
\item
  Android.mk
\end{itemize}

\begin{Shaded}
\begin{Highlighting}[]
\DataTypeTok{LOCAL\_PATH} \CharTok{:=}\StringTok{ }\CharTok{$(}\KeywordTok{call}\StringTok{ my{-}dir}\CharTok{)}

\KeywordTok{include} \CharTok{$(}\DataTypeTok{CLEAR\_VARS}\CharTok{)}

\DataTypeTok{LOCAL\_MODULE} \CharTok{:=}\StringTok{ liblive555}
\DataTypeTok{LOCAL\_CPPFLAGS} \CharTok{+=}\StringTok{ {-}fexceptions {-}DXLOCALE\_NOT\_USED=1 {-}DNULL=0 {-}DNO\_SSTREAM=1 {-}UIP\_ADD\_SOURCE\_MEMBERSHIP {-}DSOCKLEN\_T=socklen\_t}

\DataTypeTok{LOCAL\_C\_INCLUDES} \CharTok{:=}\StringTok{ }\CharTok{\textbackslash{}}
\StringTok{    }\CharTok{$(}\DataTypeTok{LOCAL\_PATH}\CharTok{)}\StringTok{ }\CharTok{\textbackslash{}}
\StringTok{    }\CharTok{$(}\DataTypeTok{LOCAL\_PATH}\CharTok{)}\StringTok{/BasicUsageEnvironment/include }\CharTok{\textbackslash{}}
\StringTok{    }\CharTok{$(}\DataTypeTok{LOCAL\_PATH}\CharTok{)}\StringTok{/UsageEnvironment/include }\CharTok{\textbackslash{}}
\StringTok{    }\CharTok{$(}\DataTypeTok{LOCAL\_PATH}\CharTok{)}\StringTok{/groupsock/include }\CharTok{\textbackslash{}}
\StringTok{    }\CharTok{$(}\DataTypeTok{LOCAL\_PATH}\CharTok{)}\StringTok{/liveMedia/include }\CharTok{\textbackslash{}}

\DataTypeTok{LOCAL\_SRC\_FILES} \CharTok{:=}\StringTok{ }\CharTok{\textbackslash{}}
\StringTok{    BasicUsageEnvironment/BasicHashTable.cpp         }\CharTok{\textbackslash{}}
\StringTok{        ...(这里把其他cpp、c文件都列到这里）}

\KeywordTok{include} \CharTok{$(}\DataTypeTok{BUILD\_SHARED\_LIBRARY}\CharTok{)}
\end{Highlighting}
\end{Shaded}

然后执行ndk-build进行编译：

\begin{Shaded}
\begin{Highlighting}[]
\ExtensionTok{ndk{-}build}\NormalTok{ NDK\_PROJECT\_PATH=. NDK\_APPLICATION\_MK=Application.mk}
\end{Highlighting}
\end{Shaded}

编译完成就可以得到libgnustl\_shared.so liblive555.so了。


\end{document}
